\documentclass[11pt]{article}
\usepackage[a4paper,margin=28mm]{geometry}
\usepackage{amsmath,amssymb,amsthm,booktabs,array,longtable}
\usepackage[T1]{fontenc}
\usepackage{lmodern}
\usepackage[colorlinks=true,linkcolor=blue,urlcolor=blue,citecolor=blue]{hyperref}
\newtheorem{theorem}{Theorem}
\newtheorem{lemma}[theorem]{Lemma}
\newtheorem{corollary}[theorem]{Corollary}
\newcommand{\tr}{\operatorname{Tr}}
\newcommand{\R}{\mathcal R}
\newcommand{\e}{e}
\newcommand{\F}{F_{k,\ell}}
\title{Extending Li--Poon--Sze's elementary-symmetric\newline
realignment comparison to unequal dimensions}
\author{Haobo Ma\\\small ChronoAI Pte Ltd, Singapore
\and Wenlin Zhang\\\small National University of Singapore\\
\small ChronoAI Pte Ltd, Singapore}
\date{November 16, 2026}
\begin{document}
\maketitle
\begin{abstract}
Extending Li--Poon--Sze's equal-dimensional answer (2011, Theorem~2.4),
we answer yes to their open problem (P2), arXiv:1103.2499v2, p.~4,
for every $2\le m<n$ and $2\le\ell\le m^2$, by combining a covariance
bound with explicit positive block states.
The maximal degree-$\ell$ elementary symmetric function of realignment
singular values over separable states is strictly smaller than its maximum
over states satisfying the realignment criterion. The original question
is Lupo--Aniello--Scardicchio (2008), Section~6, p.~15, item~(ii).
One matrix works simultaneously at every nontrivial degree, and small
white-noise mixtures give full-rank examples strictly inside the
realignment constraint.
\end{abstract}

\section{The question and the answer}
Li, Poon and Sze formulate two open problems, (P1) and (P2), in
\cite[p.~4]{LPS}. Their wording for (P2) is ``To determine if
$\widetilde B_\ell(m,n)>B_\ell(m,n)$.'' The comparison is between
separable states and states allowed by the computable cross-norm or
realignment (CCNR) criterion. They prove equality in equal dimensions
\cite[Theorem~2.4]{LPS}. We prove strict inequality in every unequal
dimension pair. The original question is item~(ii) in
\cite[Section~6, p.~15]{LAS}; its notation $(d,D)$ corresponds to
$(m^2,n^2)$ here.

A density matrix on $\mathbb C^m\otimes\mathbb C^n$ is a positive
semidefinite $mn\times mn$ matrix of trace one. It is separable if it is a
finite convex combination of product density matrices. In product bases,
define the $m^2\times n^2$ realignment by
\[
 \R(\rho)_{(i,j),(a,b)}=\rho_{(i,a),(j,b)}.
\]
Write $s_1\ge\cdots\ge s_k\ge0$ for its $k=m^2$ singular values,
including zeros, and $\|\R(\rho)\|_1=\sum_j s_j$ for its trace norm.
For a nonnegative vector $s$, let
\[
 \e_\ell(s)=\sum_{j_1<\cdots<j_\ell}s_{j_1}\cdots s_{j_\ell}.
\]
Following \cite{LPS}, write $B_\ell(m,n)$ for the maximum of this
function over separable states and $\widetilde B_\ell(m,n)$ for its
maximum over density matrices satisfying $\|\R(\rho)\|_1\le1$.
These maxima exist: both state sets are compact in finite dimensions,
and the spectral function is continuous.

For $0\le x\le1$, put $b_x=(1-x)/(k-1)$ and
\begin{equation}\label{profile}
 \F(x)=\binom{k-1}{\ell}b_x^\ell+
 x\binom{k-1}{\ell-1}b_x^{\ell-1},\qquad
 \delta=1-\sqrt{(1-1/m)(1-1/n)}.
\end{equation}
A binomial coefficient is zero when its lower index exceeds its upper
index.

\begin{theorem}\label{main}
Let $2\le m<n$, $k=m^2$, and $2\le\ell\le k$. Define
\begin{equation}\label{choose-a}
 a=\begin{cases}
 1/\sqrt{mn},&n<m^3-m,\\
 1/(m\sqrt{m^2-1}),&m^3-m\le n<m^3,\\
 1/m^2,&n\ge m^3.
 \end{cases}
\end{equation}
There is a density matrix $\rho$ with realignment singular values
$(a,b_a,\ldots,b_a)$ and trace norm one. Moreover,
\begin{equation}\label{gap}
 B_\ell(m,n)\le\F(\delta)<\F(a)
 =\e_\ell(s(\R(\rho)))\le\widetilde B_\ell(m,n).
\end{equation}
The same $\rho$ works for every $\ell=2,\ldots,m^2$.
\end{theorem}

The density-matrix family and the low- and high-dimensional constructions
are due to Li--Poon--Sze \cite[proof of Theorem~2.3]{LPS}. The covariance
estimate used below is due to Zhang, Zhang, Zhang and Guo
\cite[Theorem~1]{ZZZG}. The combination gives an explicit separable
upper bound and a strictly larger attainable value. In the middle range,
an occupied tensor block with a zero complement supplies the needed witness.
Related elementary-symmetric bounds appear in Peled, Te'eni, Carmi and
Cohen's correlation-minor framework \cite[Theorem~1]{PTCC}, with a
filter-normal-form hypothesis. The bound below applies directly to all
separable states; it differs from their proposed extension of that bound
beyond filter normal form.

Aniello and Lupo \cite[Section~4, pp.~6--8]{AL} also formulate the
question of stricter elementary-symmetric bounds independent of CCNR;
the unequal-dimensional witnesses here answer that question affirmatively.
Gittsovich, G\"uhne, Hyllus and Eisert
\cite[Proposition~IV.2 and Remark~IV.3]{GGHE} recover the covariance
estimate from their covariance-matrix criterion. Their reduced correlation
block bound (Proposition~IV.4) concerns deletion of the identity row and
column; the tail bound below instead concerns the full realignment spectrum.
The multipartite correlation-minor bounds of Lenny et al.
\cite[Theorems~1--4]{LTPCC} retain filter-normal-form or strong
filter-normal-form assumptions.

Deside, Kala and Cerf \cite[Appendix~D, equations~(D23)--(D35)]{DKC}
use elementary symmetric functions through Newton identities in a
different family of realignment-moment criteria. Their comparison is
between singular moments and weak moments of the realigned operator;
combined with the CCNR inequality it gives necessary separability
conditions. Here the comparison is between two constrained maxima, and
the witnesses satisfy CCNR while exceeding every higher-degree separable
profile bound.

\section{A separable spectral bound}
\begin{lemma}\label{tail}
For a separable state $\sigma$, its realignment singular values satisfy
\[
 \sum_{j=1}^k s_j\le1,\qquad
 \sum_{j=2}^k s_j\le\sqrt{(1-1/m)(1-1/n)}=1-\delta.
\]
\end{lemma}
\begin{proof}
Write $\sigma=\sum_h p_h A_h\otimes B_h$, where $p_h\ge0$,
$\sum_h p_h=1$, and $A_h,B_h$ are density matrices. Let $A=\sum_h p_hA_h$
and $B=\sum_h p_hB_h$ be the marginals. Use row vectorization, written
as a column $\operatorname{vec}(\cdot)$. The trace norm of
$xy^T$ is $\|x\|_2\|y\|_2$, also for complex vectors. Thus the triangle
inequality and $\tr A_h^2,\tr B_h^2\le1$ give the first bound.
Also,
\[
 \R(\sigma)-\operatorname{vec}(A)\operatorname{vec}(B)^T
 =\sum_h p_h\operatorname{vec}(A_h-A)
                  \operatorname{vec}(B_h-B)^T.
\]
Cauchy--Schwarz gives
\begin{align*}
 \|\R(\sigma-A\otimes B)\|_1
 &\le\sqrt{\Bigl(\sum_h p_h\tr A_h^2-\tr A^2\Bigr)
                   \Bigl(\sum_h p_h\tr B_h^2-\tr B^2\Bigr)}\\
 &\le\sqrt{(1-\tr A^2)(1-\tr B^2)}.
\end{align*}
This is the covariance estimate in \cite[Theorem~1]{ZZZG}.
For any rectangular matrix $C$ and rank-at-most-one matrix $D$,
$s_{j+1}(C)\le s_j(C-D)$: in the min--max formula for singular values,
intersect a test subspace with $\ker D$, losing at most one dimension.
Summing gives $\sum_{j\ge2}s_j(C)\le\|C-D\|_1$.
Apply this with $D=\operatorname{vec}(A)\operatorname{vec}(B)^T$.
Finally, the eigenvalues of a density matrix of order $m$ sum to one,
so Cauchy--Schwarz gives $\tr A^2\ge1/m$; similarly $\tr B^2\ge1/n$.
\end{proof}

\begin{lemma}\label{scalar}
If $s_1\ge\cdots\ge s_k\ge0$, $\sum_j s_j\le1$, and
$\sum_{j\ge2}s_j\le1-\delta$, where $1/k<\delta<1$, then
$\e_\ell(s)\le\F(\delta)$ for $2\le\ell\le k$.
The function $\F$ is strictly decreasing on $[1/k,1)$.
\end{lemma}
\begin{proof}
Increase the first coordinate to $x=1-\sum_{j\ge2}s_j\ge\delta$.
Every elementary symmetric function is coordinatewise nondecreasing.
For a fixed tail sum, replacing two tail coordinates $u,v$ by their mean
preserves the terms linear in $u,v$ and increases their product by
$(u-v)^2/4$. Its coefficient is the nonnegative elementary symmetric
function of degree $\ell-2$ in the remaining coordinates. Successive
averaging, and continuity, therefore bound the resulting vector by the
equal-tail value $\F(x)$. Direct differentiation yields
\[
 \F'(x)=\frac{\ell-1}{k-1}\binom{k-1}{\ell-1}
 b_x^{\ell-2}(b_x-x).
\]
This is negative for $1/k<x<1$. The endpoint $x=1$ follows by continuity.
\end{proof}

For the dimensions of Theorem~\ref{main}, $\delta<1$, and
\begin{equation}\label{delta-half}
 \delta>1-\sqrt{1-1/m}>\frac1{2m}\ge\frac1{m^2}.
\end{equation}
The middle inequality follows by squaring
$\sqrt{1-1/m}<1-1/(2m)$; both sides are nonnegative.
Lemmas~\ref{tail} and~\ref{scalar} prove the first inequality in
\eqref{gap}.

\section{Positive witnesses and their complete spectra}
Split Bob's space as $(\mathbb C^m\otimes\mathbb C^q)\oplus\mathbb C^r$.
Let $P$ and $Q$ be the projections onto these summands. For
$h=1,\ldots,q$, set
\[
 \Omega_h=\sum_{i=1}^m |i\rangle\otimes(|i\rangle\otimes|h\rangle),
 \quad E=\sum_{h=1}^q|\Omega_h\rangle\langle\Omega_h|,
 \quad \rho(t,u,v)=tE+I_m\otimes(uP+vQ).
\]
For $t,u,v\ge0$ this matrix is positive semidefinite, and
\begin{equation}\label{trace}
 \tr\rho=mq(t+mu)+mr v.
\end{equation}
Let $d=\operatorname{vec}(I_m)$, so $d^*d=m$. Writing
$C=\R(\rho)$, direct multiplication gives
\begin{equation}\label{gram}
 CC^*=\mu I_{m^2}+\nu dd^*,\qquad
 \mu=qt^2,\quad \nu=q(2tu+mu^2)+rv^2.
\end{equation}
For completeness, an occupied entry in column $((a,p),(b,h))$ is
$\delta_{ph}(t\delta_{ia}\delta_{jb}
+u\delta_{ij}\delta_{ab})$ in row $(i,j)$; an unused entry in column
$(a,b)$ is $v\delta_{ij}\delta_{ab}$, and mixed entries vanish.
The $tt$, $tu$, $uu$, and $vv$ products are respectively
$qt^2I$, $2qtu\,dd^*$, $qmu^2dd^*$, and $rv^2dd^*$,
which proves \eqref{gram}. Its eigenvalues are
\begin{equation}\label{eigen}
 \mu\quad(k-1\text{ times}),\qquad
 \mu+m\nu=q(t+mu)^2+mr v^2\quad(\text{once}).
\end{equation}

\paragraph{The range $m<n<m^3-m$.}
Write $n=mq+r$, $0\le r<m$, and put
\[
 a=\frac1{\sqrt{mn}},\quad b=\frac{1-a}{m^2-1},\quad
 t=\frac b{\sqrt q},\quad u=a^2-\frac b{m\sqrt q},\quad v=a^2.
\]
Here $1\le q\le m^2-2$. We check $u\ge0$ rather than assuming
positivity from the desired spectrum. Let $K=m^2-2$. The convex quadratic
$g(z)=[m(z+1)-1]^2-m^4z$ is nonpositive for $1\le z\le K$:
$g(1)\le0$ because $2m-1\le m^2$, while
\[
 -g(K)=2m^3-m^2-2m-1\ge3m^2-2m-1>0 \quad(m\ge2).
\]
Since $n\le m(q+1)-1$, this proves $n\le m^2\sqrt q$.
Consequently
\[
 \frac n{\sqrt q}-\sqrt{\frac n{mq}}\le m^2-1.
\]
Multiplying out the definition of $u$ shows that this inequality is exactly
$u\ge0$. Also $t,v\ge0$ and $t+mu=ma^2$.
Equations~\eqref{trace}--\eqref{eigen}, with $mn a^2=1$, now give
$\tr\rho=1$, $\mu=b^2$, and $\mu+m\nu=a^2$.

\paragraph{Supported blocks.}
More generally, suppose $1\le q\le m^2$ and $mq\le n$.
Take $r=n-mq$ and use
\[
 a=\frac1{m\sqrt q},\quad b=\frac{1-a}{m^2-1},\quad
 t=\frac b{\sqrt q},\quad u=a^2-\frac b{m\sqrt q},\quad v=0.
\]
The condition $u\ge0$ is equivalent to $b\le a$, which follows from
$\sqrt q\le m$. Again \eqref{trace}--\eqref{eigen} give trace one and
squared singular values $a^2,b^2,\ldots,b^2$, using $m^2q a^2=1$.
In the middle range of \eqref{choose-a}, take $q=m^2-1$; in the last
range take $q=m^2$. The latter gives $a=b=1/m^2$, $u=v=0$, and
$t=1/m^3$. In all ranges the trace norm is $a+(m^2-1)b=1$.

\paragraph{The strict comparison.}
In the first range, $a\ge1/m^2$, and $a<\delta$ follows from
\[
 (\sqrt{mn}-1)^2-(m-1)(n-1)=(\sqrt n-\sqrt m)^2>0.
\]
Indeed $\sqrt{mn}-1>0$, so this is equivalent to
$\sqrt{(m-1)(n-1)}<\sqrt{mn}-1$.
For the supported middle range, $a>1/m^2$. If $m\ge3$, then
$\sqrt{m^2-1}>2$, so $a<1/(2m)<\delta$ by \eqref{delta-half}.
For $m=2$, the exact comparisons
\[
 \frac1{2\sqrt3}<\frac7{24}<1-\frac1{\sqrt2}<\delta
\]
follow by squaring positive quantities: $144<147$ and $288<289$.
For the last range, $a=1/m^2<\delta$ already follows from
\eqref{delta-half}. Thus $1/k\le a<\delta<1$ in all three cases.
Lemma~\ref{scalar} gives $\F(\delta)<\F(a)$, completing the proof of
Theorem~\ref{main}.

\begin{corollary}\label{entangled}
For each $2\le m<n$ there is an entangled state satisfying CCNR whose
degree-$\ell$ realignment elementary symmetric function exceeds the
separable maximum simultaneously for all $\ell=2,\ldots,m^2$.
For degree one, $B_1(m,n)=\widetilde B_1(m,n)=1$.
\end{corollary}
\begin{proof}
Use Theorem~\ref{main}; separability would contradict any one of its strict
inequalities. For degree one CCNR gives the upper bound one, attained by
a pure product state.
\end{proof}

\section{Sibling questions and consequences}
The answer to (P2) has a complete dimension dichotomy. Some associated
extremal formulas were already obtained in the posing papers; the following
statements separate those formulas from the unequal-dimensional separation.

\begin{corollary}\label{classification}
For $m,n\ge2$ and $2\le\ell\le\min(m,n)^2$, the answer to (P2) is
yes exactly when $m\ne n$. If $m=n$, then
\[
 B_\ell(m,m)=\widetilde B_\ell(m,m)
 =\F(1/m),\qquad k=m^2.
\]
In particular, the exact values at $m=n=2$ and degrees $1,2,3,4$ are
$1,1/3,5/108,1/432$.
\end{corollary}
\begin{proof}
Swapping the two tensor factors transposes the realignment, up to
permutations, and hence preserves its singular values. Thus
Theorem~\ref{main} covers both unequal orders.
For equal dimensions, the tail bound has $\delta=1/m$ and gives
$B_\ell\le\F(1/m)$. For every density matrix, testing its realignment
against $\operatorname{vec}(I_m)/\sqrt m$ and
$\operatorname{vec}(I_n)/\sqrt n$ gives $s_1\ge1/\sqrt{mn}$.
Consequently, in equal dimensions, CCNR gives
$\sum_{j\ge2}s_j\le1-1/m$, and Lemma~\ref{scalar} bounds
$\widetilde B_\ell$ by the same value.
The separable state
\[
 \frac{I_{m^2}+|\Omega\rangle\langle\Omega|}{m(m+1)},
 \qquad \Omega=\sum_{i=1}^m|ii\rangle,
\]
attains it: it is the uniform unit-vector average of
$|u\rangle\langle u|\otimes|\bar u\rangle\langle\bar u|$, and its
realignment spectrum is $1/m,1/(m(m+1)),\ldots,1/(m(m+1))$.
The averaging identity follows from
$\mathbb E[u_i\bar u_j\bar u_a u_b]
=(\delta_{ij}\delta_{ab}+\delta_{ia}\delta_{jb})/(m(m+1))$.
An average belongs to the compact convex hull of product states.
This recovers Li--Poon--Sze's Theorem~2.4, including the two-qubit values
estimated in \cite[Table~1]{LAS}. Degree one is Corollary~\ref{entangled}.
\end{proof}

\begin{corollary}\label{naive}
Let $2\le m\le n$ and $2\le\ell\le k=m^2$. The naive CCNR upper
bound $\binom{k}{\ell}k^{-\ell}$ is attained if and only if $n\ge m^3$.
For $n<m^3$, writing $\alpha=1/\sqrt{mn}$ gives
\[
 \widetilde B_\ell(m,n)\le\F(\alpha)
 <\binom{k}{\ell}k^{-\ell}.
\]
\end{corollary}
\begin{proof}
Averaging all coordinates under the constraint $\sum s_j\le1$ gives
the naive bound. If $n<m^3$, then $s_1\ge\alpha>1/k$ as above,
so the tail sum is at most $1-\alpha$. Apply Lemma~\ref{scalar};
strict monotonicity compares $\alpha$ to $1/k$.
If $n\ge m^3$, the uniform witness in the proof of
Theorem~\ref{main} attains the bound.
The strict side recovers \cite[Proposition~4]{LAS}, whose condition
$D<d^3$ becomes $n<m^3$; the attained side is
\cite[Theorem~2.3]{LPS}.
\end{proof}

For the exact-maxima question (P1), Li--Poon--Sze already prove
\cite[Theorem~2.3]{LPS}
\begin{equation}\label{known-maxima}
 \widetilde B_\ell(m,n)=
 \begin{cases}
 \F(1/\sqrt{mn}),&m\le n\le m^3-m/2,\\
 \binom{k}{\ell}k^{-\ell},&n\ge m^3.
 \end{cases}
\end{equation}
In particular, at $(m,n)=(2,3)$ our bound gives
$B_\ell(2,3)\le F_{4,\ell}(1-1/\sqrt3)
<F_{4,\ell}(1/\sqrt6)=\widetilde B_\ell(2,3)$.
This proves the strict comparison suggested by \cite[Table~2]{LAS},
without assigning exact values to that table's numerical separable maxima.
For the exceptional case highlighted in \cite[p.~9]{LPS},
the supported witness gives the explicit bracket
\[
 B_\ell(3,26)\le F_{9,\ell}(1-5/\sqrt{39})
 <F_{9,\ell}(1/\sqrt{72})
 \le\widetilde B_\ell(3,26)\le F_{9,\ell}(1/\sqrt{78}).
\]
The exact CCNR maximum in $m^3-m/2<n<m^3$, and exact unequal-dimensional
separable maxima, are not determined by these bounds.

\begin{corollary}\label{interior}
In every unequal dimension pair there are positive definite density matrices
with $\|\R(\rho)\|_1<1$ and
$\e_\ell(s(\R(\rho)))>B_\ell(m,n)$ simultaneously for all
$2\le\ell\le\min(m,n)^2$.
More precisely, for $m<n$, take the witness of Theorem~\ref{main} and put
\[
 \varepsilon_0=\frac12\min_{2\le\ell\le k}
 \frac{1- F_{k,\ell}(\delta)/F_{k,\ell}(a)}{\ell}>0.
\]
Every $\rho_\varepsilon=(1-\varepsilon)\rho+
\varepsilon I_{mn}/(mn)$ with $0<\varepsilon\le\varepsilon_0$
has these properties.
\end{corollary}
\begin{proof}
The least eigenvalue is at least $\varepsilon/(mn)>0$, and convexity gives
\[
 \|\R(\rho_\varepsilon)\|_1
 \le1-\varepsilon+\varepsilon/\sqrt{mn}<1.
\]
Substitute $t_\varepsilon=(1-\varepsilon)t$,
$u_\varepsilon=(1-\varepsilon)u+\varepsilon/(mn)$, and
$v_\varepsilon=(1-\varepsilon)v+\varepsilon/(mn)$ in \eqref{gram}.
Using $t+mu=ma^2$, the realignment singular values are
\[
 a_\varepsilon=
 \sqrt{(1-\varepsilon)^2a^2+(2\varepsilon-\varepsilon^2)/(mn)},
 \qquad b_\varepsilon=(1-\varepsilon)b_a
 \quad(k-1\text{ times}).
\]
For the main construction the formula uses $mn a^2=1$; for the supported
construction it uses $r v=0$ and $m^2q a^2=1$.
Thus $a_\varepsilon\ge(1-\varepsilon)a$, so coordinatewise monotonicity
and homogeneity give
\[
 \e_\ell(s(\R(\rho_\varepsilon)))
 \ge(1-\varepsilon)^\ell F_{k,\ell}(a)
 \ge(1-\ell\varepsilon)F_{k,\ell}(a)>F_{k,\ell}(\delta).
\]
The last inequality follows from the definition of $\varepsilon_0$;
the middle one is Bernoulli's inequality. Apply Theorem~\ref{main}.
\end{proof}

The same states give a channel consequence of
\cite[Section~5, equations~(38)--(43)]{LAS}. A channel is a completely
positive trace-preserving linear map between matrix spaces; it is
entanglement breaking if its action on one part of any bipartite state
always gives a separable state.

\begin{corollary}\label{channels}
For every $2\le m<n$ there is a channel $\Phi:M_m(\mathbb C)\to
M_n(\mathbb C)$ which is not entanglement breaking, although its
Hilbert--Schmidt matrix has singular-value sum at most $m$.
At every degree $2\le\ell\le m^2$, its elementary symmetric value
exceeds the upper bound $m^\ell F_{m^2,\ell}(\delta)$ valid for
entanglement-breaking channels. The sum can be made strictly smaller
than $m$, with a positive definite normalized Choi matrix.
\end{corollary}
\begin{proof}
For each block state above,
$\operatorname{Tr}_B\rho=[q(t+mu)+rv]I_m=I_m/m$.
Thus $\Phi(X)=m\operatorname{Tr}_A[(X^T\otimes I_n)\rho]$
is a channel with normalized Choi matrix $\rho$.
Its Hilbert--Schmidt matrix is $m\R(\rho)^T$, so its singular values
are $m$ times those of $\R(\rho)$.
An entanglement-breaking channel has separable normalized Choi matrix;
Lemma~\ref{tail} and Lemma~\ref{scalar} give the stated upper bound.
Theorem~\ref{main} and Corollary~\ref{interior} give the strict
comparisons and the positive definite variant.
\end{proof}

\section{Coverage and mathematical evidence}
{\small
\begin{longtable}{@{}p{.43\linewidth}p{.53\linewidth}@{}}
\toprule Source item & Coverage\\\midrule\endhead
Li--Poon--Sze, p.~4 (P2), unequal dimensions &
Answered: yes. Theorem~\ref{main}, in both tensor orders; one witness
for all nontrivial degrees.\\
Li--Poon--Sze (P2), equal dimensions & Answered: no.
Reproved; prior proof Li--Poon--Sze (2011), Theorem~2.4.
Corollary~\ref{classification}.\\
Lupo--Aniello--Scardicchio, Section~6(ii), p.~15 &
Answered: yes in unequal dimensions; no in equal dimensions.\\
Aniello--Lupo, Section~4, p.~8, independent stricter bounds &
Answered: yes in unequal dimensions, by Theorem~\ref{main}; the
equal-dimensional obstruction is Corollary~\ref{classification}.\\
Lupo--Aniello--Scardicchio, Section~8, p.~18, conjecture on existence
of strengthened independent criteria & Proved in every unequal dimension
pair, by Theorem~\ref{main}.\\
Lupo--Aniello--Scardicchio, Proposition~4, naive upper bound &
Reproved; prior proof Lupo--Aniello--Scardicchio (2008) for $n<m^3$;
Li--Poon--Sze (2011) for $n\ge m^3$. Corollary~\ref{naive}.\\
Lupo--Aniello--Scardicchio, Tables~1--2 & Table~1 values:
reproved, prior proof Li--Poon--Sze (2011). Table~2 strict comparison:
proved; exact separable values not covered.\\
Li--Poon--Sze (P1); Lupo--Aniello--Scardicchio Section~6(i) &
Prior proof cited for \eqref{known-maxima}; exact remaining CCNR maxima
and unequal-dimensional separable maxima not covered.\\
Li--Poon--Sze, p.~9, $(3,26)$ & Explicit bracket above;
the exact maximum not covered.\\
Peled et al., Theorem~1 and its proposed unrestricted extension &
FNF bound: prior proof cited. Removing FNF from that stronger bound:
not covered.\\
Full rank and strict CCNR inequality & Not posed in the primary source;
proved, Corollary~\ref{interior}.\\
Lupo--Aniello--Scardicchio, Section~5, channel formulation &
Not posed as an existence question there; proved consequence,
Corollary~\ref{channels}.\\\bottomrule
\end{longtable}}

Theorem~\ref{main} and the corollaries have written proofs above;
\eqref{known-maxima} is a cited prior result.
The accompanying \texttt{verify.py} checks exact trace and realignment
row-Gram identities, polynomial identities, and rational interval comparisons
for the listed block-state and white-noise examples. Its finite ranges and
interval precision are specified in \texttt{verification.json}.

At the pinned source revision, the declaration
\href{https://github.com/the-omega-institute/trureturing/blob/d13f56f01c2e4710683d27a74c67cba9ec344e19/D5/S3/Quantum/Entanglement/Witness/LiPoonSzeP2.lean#L806}{\texttt{LiPoonSzeP2.result}}
states, for each $m,n,\ell$ in the theorem's ranges, the existence of an
actual complex density matrix satisfying CCNR and exceeding a common upper
bound for every separable state's degree-$\ell$ value.
Its quantifiers are $\forall m,n,\ell\;\exists\rho$.
The same file's \texttt{main\_witness} ($n<m^3-m$),
\texttt{supported\_witness} ($1\le q\le m^2$, $mq\le n$), and
\texttt{uniform\_witness} ($n\ge m^3$) specify the three spectra
independently of $\ell$; the written argument assembles one simultaneous
witness. The definitions of \texttt{realign}, \texttt{singularVector},
\texttt{esymmSV}, and the separable profile bound are in
\href{https://github.com/the-omega-institute/trureturing/blob/d13f56f01c2e4710683d27a74c67cba9ec344e19/D5/S3/Quantum/Entanglement/Witness/RealignmentSchurBound.lean#L924}{\texttt{RealignmentSchurBound.lean}}:
\texttt{separable\_esymm\_bound} gives $\e_\ell\le F_{k,\ell}(\delta)$.
The covariance and marginal-purity tail inequalities in
\href{https://github.com/the-omega-institute/trureturing/blob/d13f56f01c2e4710683d27a74c67cba9ec344e19/D5/S3/Quantum/Entanglement/Witness/CorrelationMinorNormCovariance.lean#L162}{\texttt{CorrelationMinorNormCovariance.lean}}
are \texttt{actual\_covariance\_nuclear} and
\texttt{correlation\_tail\_purity}; they concern correlation matrices
in orthonormal Hermitian bases, which have the realignment singular values.
The full-rank white-noise extension, channel consequence and
equal-dimensional averaging argument are written proofs here, separate
from those declarations.

\begin{thebibliography}{9}
\bibitem{LAS} C.~Lupo, P.~Aniello and A.~Scardicchio,
\emph{Bipartite quantum systems: on the realignment criterion and beyond},
J.~Phys.~A: Math.~Theor. \textbf{41} (2008), 415301.
\href{https://doi.org/10.1088/1751-8113/41/41/415301}{doi:10.1088/1751-8113/41/41/415301};
\href{https://arxiv.org/abs/0802.2019v2}{arXiv:0802.2019v2}.
\bibitem{LPS} C.-K.~Li, Y.-T.~Poon and N.-S.~Sze,
\emph{A note on the realignment criterion},
J.~Phys.~A: Math.~Theor. \textbf{44} (2011), 315304.
\href{https://doi.org/10.1088/1751-8113/44/31/315304}{doi:10.1088/1751-8113/44/31/315304};
\href{https://arxiv.org/abs/1103.2499v2}{arXiv:1103.2499v2}.
\bibitem{ZZZG} C.-J.~Zhang, Y.-S.~Zhang, S.~Zhang and G.-C.~Guo,
\emph{Entanglement detection beyond the computable cross-norm or
realignment criterion}, Phys.~Rev.~A \textbf{77} (2008), 060301(R).
\href{https://doi.org/10.1103/PhysRevA.77.060301}{doi:10.1103/PhysRevA.77.060301};
\href{https://arxiv.org/abs/0709.3766v6}{arXiv:0709.3766v6}.
\bibitem{PTCC} B.~Y.~Peled, A.~Te'eni, A.~Carmi and E.~Cohen,
\emph{Correlation minor norms, entanglement detection and discord},
Scientific Reports \textbf{11} (2021), 2849.
\href{https://doi.org/10.1038/s41598-021-82303-3}{doi:10.1038/s41598-021-82303-3};
\href{https://arxiv.org/abs/2005.12079}{arXiv:2005.12079}.
\bibitem{AL} P.~Aniello and C.~Lupo,
\emph{On the relation between Schmidt coefficients and entanglement},
Open Syst.~Inf.~Dyn. \textbf{16} (2009), 127--143.
\href{https://doi.org/10.1142/S1230161209000104}{doi:10.1142/S1230161209000104};
\href{https://arxiv.org/abs/0812.4167v1}{arXiv:0812.4167v1}.
\bibitem{GGHE} O.~Gittsovich, O.~G\"uhne, P.~Hyllus and J.~Eisert,
\emph{Unifying several separability conditions using the covariance
matrix criterion}, Phys.~Rev.~A \textbf{78} (2008), 052319.
\href{https://doi.org/10.1103/PhysRevA.78.052319}{doi:10.1103/PhysRevA.78.052319};
\href{https://arxiv.org/abs/0803.0757v3}{arXiv:0803.0757v3}.
\bibitem{LTPCC} R.~Lenny, A.~Te'eni, B.~Y.~Peled, A.~Carmi and E.~Cohen,
\emph{Multipartite entanglement detection via correlation minor norm},
Quantum Inf.~Process. \textbf{22} (2023), 292.
\href{https://doi.org/10.1007/s11128-023-04046-6}{doi:10.1007/s11128-023-04046-6};
\href{https://arxiv.org/abs/2302.00650v1}{arXiv:2302.00650v1}.
\bibitem{DKC} S.~Deside, V.~Kala and N.~J.~Cerf,
\emph{Entanglement detection from moments of the realigned state:
application to continuous-variable quantum systems},
\href{https://arxiv.org/abs/2611.17021v1}{arXiv:2611.17021v1} (2026).
\end{thebibliography}
\medskip\noindent\small
\copyright\ 2026 Haobo Ma and Wenlin Zhang.
Licensed under Creative Commons Attribution 4.0 International (CC BY 4.0).
\end{document}
